\section{A genuine normal graph contact direction and whole coefficient slice}

The marked ideal remains on its full smooth ambient space. The arbitrary
singular graph is never substituted for its smooth marked ambient stratum.
Canonical invariant tails, maximal centres and history descent remain separate.

\begin{theorem}[A literal graph-normal equation is an actual marking-one tangent direction]
\label{hir:actual-graph-normal-tangent-direction}
\uses{hir:actual-graph-functional-ideal-every-point,mc:tangent-direction,bcalc:hypersurface}
\lean{Hironaka.HironakaActualGraphContactCoefficient.graphFunction_chartDeriv_eq_one,Hironaka.HironakaActualGraphContactCoefficient.graphFunction_isTangentDirection,Hironaka.HironakaActualGraphContactCoefficient.graphNormalContact_regularAtDim}
\leanok
On any actual open where $f$ is holomorphic, each literal normal graph
equation $u_i=y_i-f_i(x)$ is a marking-one tangent direction of the
independent original graph vanishing ideal on the corresponding true
ambient open. Its actual preferred-chart normal derivative is exactly one.
Its actual zero hypersurface on that true open is regular of dimension
one less than the smooth ambient dimension.
\end{theorem}
\begin{proof}
The preferred chart on the actual vector model is the identity. The actual
complex chain rule gives $Du_i$ as the $i$th normal projection minus the
$i$th component of $Df$ composed with the first block projection. On the
$i$th normal basis vector these terms are one and zero. Local holomorphy
of $f$ proves holomorphy of $u_i$ on the true ambient open. Every original
graph point satisfies $u_i=0$ literally. Thus $u_i$ belongs to the entire
original graph vanishing ideal at every point of the true open. At mark
one its derivative-zero iterate is that original ideal, so these facts
give every actual tangent-direction requirement, including the nonzero
derivative at every zero, without a supplied contact certificate. The genuine
hypersurface inverse-function theorem proves actual regularity of this zero
hypersurface on the same true open.
\end{proof}

\begin{theorem}[The genuine whole graph-section ideal pullback]
\label{hir:actual-graph-section-whole-ideal-pullback}
\uses{hir:actual-graph-functional-ideal-every-point}
\lean{Hironaka.HironakaActualGraphContactCoefficient.graphSection_pullback_mem_iff}
\leanok
For every original subset $W$, point $x$, and $f$ holomorphic only near
$x$, the entire independently defined functional pullback of the graph
vanishing ideal along the literal graph section at $x$ is the original
vanishing ideal of $W$, for every original scalar numerator. No coherence
or pullback equality is input.
\end{theorem}
\begin{proof}
For each original graph-vanishing function, true continuity of the actual
holomorphic graph section transports its entire neighborhood vanishing
relation to $W$ near $x$, and the actual complex chain rule proves
holomorphy of its composition. Every pullback generator is therefore a
base-vanishing germ. Any true finite holomorphic span remains base
vanishing: intersect the finitely many actual generator neighborhoods
with the genuine span-relation neighborhood and evaluate the finite sum.
Conversely, an original base-vanishing function $a$ gives the actual
ambient function $a\circ p$. True projection continuity transports its
whole vanishing relation to the original graph, and the ordinary complex
chain rule gives its local holomorphy. Its pullback through the actual
graph section is exactly $a$ because $p\circ\operatorname{graph}(f)=1$.
The actual one-term span proves the reverse entire membership.
\end{proof}

\begin{theorem}[The actual marking-one coefficient slice removes one genuine graph-normal equation]
\label{hir:actual-graph-marking-one-coefficient-slice}
\uses{hir:actual-graph-section-whole-ideal-pullback,hir:actual-graph-normal-tangent-direction,mc:coefficient-ideal}
\lean{Hironaka.HironakaActualGraphContactCoefficient.coefficientGraphSection_mem_iff,Hironaka.HironakaActualGraphContactCoefficient.exists_scalarGraphContactCoefficientSlice}
\leanok
Mark the original graph ideal by one, on its full smooth ambient space
with empty boundary. Its actual coefficient ideal has entire graph-section
pullback equal to the original base ideal at every point where $f$ is
holomorphic nearby. In particular let the base subset be any actual
lower graph and let $f$ have one scalar graph-normal component. The smooth
contact parametrization removes exactly that added normal equation and
recovers the whole original lower graph ideal. The actual point lies in an
internally chosen true open; its normal equation is an actual marking-one
tangent direction there, and its contact hypersurface is genuinely regular
of the lower ambient dimension. No smoothness of the lower
graph or coherence certificate is assumed for this membership dictionary.
\end{theorem}
\begin{proof}
The genuine marking-one coefficient construction contains only the
derivative-zero summand, with exponent one and mark one. Expanding the
actual finite-span and ideal-power definitions, every generator and each
finite holomorphic combination belongs to the original graph vanishing
ideal, by its actual vanishing relation. Conversely any original
vanishing germ occurs as a one-term combination with the constant unit.
Thus the coefficient ideal is that entire original graph ideal; this is
proved from true closure of vanishing germs rather than inferred from
coherence. Apply the preceding whole graph-section pullback theorem.
For the one-equation assertion substitute the actual lower graph for
$W$ and the genuine scalar normal function for $f$. The target graph is
literally the lower graph with one added graph equation; its contact
parametrization and projection are the actual graph section and block
projection. Choose a true original holomorphic neighborhood from the original
near-point holomorphy. The preceding derivative-one and actual hypersurface
regularity theorems supply this same genuine smooth graph-normal contact
direction and its actual dimension, while the whole ideal equivalence holds
for every original scalar numerator. No canonical tail or maximal-centre
invariance is inferred from the coefficient calculation.
\end{proof}
